\documentclass[a4paper,12pt,twoside,openright]{memoir}

%------------------------------------------------------
%   PACKAGES SET-UP
%------------------------------------------------------

%---------------
%   PAGE
%---------------

\usepackage{geometry}   % Required for adjusting page dimensions and margins

\geometry{
	paper=a4paper,      % Using standard european paper size, 'letterpaper' for US letter size
	top=18mm,           % Top margin (it was 18mm before, in case I fucked this up)
	bottom=12.5mm,        % Bottom margin (it was 13mm before, in case I fucked this up) % canonically changed to 12.5
%	left=21mm,          % Left margin
%	right=21mm,         % Right margin
    inner=19mm,         % inner margin
    outer=21mm,         % outer margin
	headheight=19mm,    % Header height
	%includehead,       % -- Without these, the header and footer are --
	%includefoot,       % -- not included within the area defined by margins --
	%footskip=14mm,     % Space from the bottom margin to the baseline of the footer
	%headsep=10mm,      % Space from the top margin to the baseline of the header
	%showframe,         % Uncomment to show how the type block is set on the page
}

\usepackage{xcolor}     % Required for specifying colors by name
\definecolor{mainoffwhite}{RGB}     {244, 239, 230}     % Main colour for pages and some text
\definecolor{corneryellow}{RGB}     {250, 183, 84 }     % Corners, some important things
\definecolor{purplevivid}{RGB}      {162, 92 , 156}     % Inserts
\definecolor{purplight}{RGB}        {180, 149, 184}     % Inserts highlights
\definecolor{somegrey}{RGB}         {106, 109, 114}     % Some stuff's backdrop
\definecolor{almostblack}{RGB}      {38 , 37 , 44 }     % main text colour
\definecolor{prettygreen}{RGB}      {101, 194, 149}     % A pretty green for highlighting text
\definecolor{chalgreen}{RGB}        {142, 211, 124}     % Background color for Challenge page
\definecolor{watermelony}{RGB}      {234, 80 , 151}     % Watermelony Challenge color

\usepackage[pagecolor={mainoffwhite}, nopagecolor={mainoffwhite}]{pagecolor} % provides \thepagecolor, better usage for colour consistency

\usepackage{fancyhdr}   % Fancy headers and footers!
\pagestyle{fancy}       % Set page style to use fancyhdr
\fancyhf{}

\renewcommand\chaptermark[1]{\markboth{\headingfont #1}{}}

\newsavebox{\myheadbox }% Heading storage box for centering headers in yellowcorner


\renewcommand\headrulewidth{0pt} % no line between document and header
\fancyhead{} % clear header
%\fancyhead[RE]{\nouppercase\headingfont\bfseries\fontsize{16pt}{16pt} \leftmark}
\fancyhead[RE]{\savebox{\myheadbox}{\hspace{75mm}}\makebox[\wd\myheadbox][c]{\nouppercase\headingfont\bfseries\fontsize{16pt}{16pt} \leftmark}} %\usebox{\myheadbox}
\fancyfoot{} % clear footer
\fancyfoot[RE,LO]{\headingfont\bfseries\fontsize{10pt}{10pt} Page \thepage}

\usepackage{CharisSIL}

\usepackage{float}                              % manipulate float placement
\usepackage{tikz, tikz-qtree}                   % Required to draw custom shapes & trees
\usepackage{Required/tikzvowel}                 % Vowel charts w/ tikz
\usepackage{pgf}
\usetikzlibrary{positioning, arrows.meta}
\usetikzlibrary{backgrounds, matrix, chains, positioning, decorations.pathreplacing, arrows}
\usetikzlibrary{shapes, arrows, positioning, calc, chains, scopes}
\usepackage{afterpage}                          % For changing page colors (for cover designs)

% needed for Baerian Notation (https://github.com/LLBlumire/recursive-baerian-phonotactics-notation/blob/master/Recursive_Baerian_Syntax_Notation.pdf)
%\usepackage{amsmath}
\usepackage{mathtools}
\usepackage{mathabx}
\usepackage{relsize}
\usepackage[nogroupskip,nonumberlist]{glossaries}
\usepackage{trimspaces}
\usepackage{glossary-mcols}

\usepackage{multicol}

\usepackage{array}                                          % Enables the following,
\newcolumntype{$}{>{\global\let\currentrowstyle\relax}}     % which is really just some automated formatting
\newcolumntype{^}{>{\currentrowstyle}}                      % for tables to have their headers in bold
\newcommand{\rowstyle}[1]{\gdef\currentrowstyle{#1}%
  #1\ignorespaces
}

% Required for including pictures
\usepackage{graphicx}
\graphicspath{ {./Pictures/} }
\usepackage{background}             % Required to draw and plac eiamges in background

\usepackage{ifthen}                 % Conditions!

\usepackage{lingmacros}     % Various useful things some may be used to when writing for conlanging
\usepackage{enumitem}       % Better lists
\setlist{noitemsep,nosep}   % Reduce spacing between bullet points and numbered lists, and between the items

\usepackage{lipsum}         % Inserts dummy text

\usepackage{titlesec}       % Allows to style the titles of sections and all
\usepackage{titletoc}       % Allows manipulation of the ToC

% baarux package and setup from u/akamchinjir
\usepackage[mcolblock,nostandards]{Required/baabbrevs}
\usepackage[canonical,nochapterlabels]{Required/baarux-0.9.11}
\baaruset{junctureoutersep=0.4ex, junctureinnersep=0.2ex}
\baaruset{hangindent=1.5em} % adjusting indenting on multiline glosses, specifically for HidingWaters
\renewcommand{\glossarypreamble}{\thispagestyle{empty}} % for removing page number from abbreviations glossary

\newbaarulinetype{tr}{lit}{\textit{lit.} \baarudquote}
\newbaarulinetype{tr}{alttr}{Or: \baarudquote}

% sets example counter to only count within chapters, but not show the chapter number
%\counterwithout{exesi}{chapter} % Removed this, added [nochapterlabels] to package options instead

\usepackage{verbatimbox}

\usepackage{tabu} % Tables!

\usepackage[hidelinks,hypertexnames=false]{hyperref} % Links will not alter the format

\usepackage{bookmark} % Allows for pdfbookmark

\usepackage[scaled=.92]{helvet}%. Sans serif - Helvetica
\usepackage{calc, color}

\usepackage{multirow} % For annoying tables

\usepackage{ulem}   % for better, customisable underlines


%------------------------------------------------------
%   FRILL & EMBELLISHMENTS
%------------------------------------------------------

\newcommand{\separ}{
	\begin{center}
		\noindent\rule{0.7\textwidth}{0.4pt}
	\end{center}
}

\newfontfamily\playfair[Path=Required/Fonts/]{PlayfairDisplay.ttf}
\DeclareTextFontCommand{\pffont}{\playfair}

\newfontfamily\latolight[Path=Required/Fonts/]{Lato-Light.ttf}
\DeclareTextFontCommand{\llfont}{\latolight}

\newfontfamily\nototc[Path=Required/Fonts/]{NotoSerifTC-Medium.otf}
\DeclareTextFontCommand{\tcfont}{\nototc}

\newfontfamily\karma[Path=Required/Fonts/]{Karma-Medium.ttf}
\DeclareTextFontCommand{\kmfont}{\karma}

\usepackage{xeCJK} % for East Asian scripts, must be loaded with XeLaTeX

\usepackage{tabto}

\NumTabs{5}

\newcommand{\scell}[2][c]{%
  \begin{tabular}[#1]{@{}c@{}}#2\end{tabular}} % Used for line breaks within a cell; [t] aligns with top, [b] aligns with bottom

%------------------------------------------------------
%   TABLE OF CONTENTS & BAABBREVS GLOSSARY
%------------------------------------------------------
\usepackage{pdfpages} % for inserting full-page pdf?

\usepackage{tocloft}

\makeglossaries

\makeatletter
\renewcommand{\contentsname}{Showcases}
\renewcommand{\cftchapteraftersnum}{ | }
\renewcommand{\cftchapterfont}{\headingfont}
\renewcommand{\cftchapterleader}{\bfseries\cftdotfill{\cftdotsep}}
\renewcommand{\cftchapterpagefont}{\headingfont}
\makeatother

\newlength{\mylen}   % a "scratch" length
\settowidth{\mylen}{\cftchapterpresnum\cftchapteraftersnum} % extra space
\addtolength{\cftchapternumwidth}{\mylen} % add the extra space


%------------------------------------------------------
%   TITLE STYLING
%------------------------------------------------------

% \newfontfamily\headingfont[]{Arial} % creates a new font command we'll use for titles and headings
% \newfontfamily\latofont[]{Lato}     % creates a new font command we'll use for author info
\newfontfamily\headingfont[]{NotoSans-Regular} % creates a new font command we'll use for titles and headings
\newfontfamily\latofont[]{NotoSans-Regular}    % creates a new font command we'll use for author info
\newcommand{\titlesize}{\fontsize{32pt}{32pt}\selectfont}   % creates a command to set font size to 32pt
\newcommand{\authorsize}{\fontsize{16pt}{16pt}\selectfont}  % creates a command to set font size to 16pt



\makeatletter
\def\@maketitle{
  %\newpage
  \null
  %\vskip .5em
  \begin{center}
  %\let \footnote \thanks
    {\headingfont\titlesize\bfseries \@title \par}
    \vskip .5em
    {%\large
      %\lineskip .5em
      \begin{tabular}[t]{c}
        {\latofont\authorsize\bfseries by \textit{\@author}}
      \end{tabular}\par}
    \vskip 1em
    %{\large \@date}
    %\vskip 1em
    {\headingfont\authorsize\bfseries \subtitle}
  \end{center}
  \par
  \vskip 1.5em}
\makeatother

%---------------
%   TITLE
%---------------


%---------------
%   SUBTITLE
%---------------



%---------------
%   AUTHOR
%---------------


%------------------------------------------------------
%   GRAPHICS
%------------------------------------------------------



%---------------
%   HEADER
%---------------

\renewcommand{\headrulewidth}{0.0pt}

% Defining shapes with tikz
\tikzset{
    hdrmain/.style={ % requires library shapes.geometric
        draw,
        fill=corneryellow,
%        inner sep=mm,
%        outer sep=0mm,
        trapezium,
        trapezium left angle=145,
        trapezium right angle=90,
        minimum height=13mm,
        minimum width= 123mm,
%        text width=45mm,
        text centered,
    },
}
%        \begin{tikzpicture}
%            \node[hdrmain](hdr){\color{almostblack}Plenty of stuff};
%        \end{tikzpicture}

\newif\ifBgUse % Lys's test for removing bg from pages

\backgroundsetup{scale=1,angle=0,opacity=1,contents={}}
\AddEverypageHook{%
    \ifBgUse % Lys's test for removing bg from pages
    \backgroundsetup{contents={}
}%  Set an if/else conditional statement that checks whether the page is odd or not
    \ifthenelse{\isodd{\value{page}}}
        {% IF YES
            \SetBgPosition{current page.north west}
            \SetBgContents{
                \begin{tikzpicture}[remember picture,overlay,xshift=0mm,yshift=0mm]
                    \fill[color=corneryellow](0mm,-297mm) -- (0mm,-290.5mm) -- (95.8mm,-290.5mm) -- (105mm,-297mm) -- cycle;
                \end{tikzpicture}
            }
            \BgMaterial%
        }%
        {% IF NOT
            \SetBgPosition{current page.north east}
            \SetBgContents{
                \begin{tikzpicture}[remember picture,overlay,xshift=0mm,yshift=0mm]
                    \fill[color=corneryellow](0mm,0mm) -- (0mm,-13mm) -- (-105mm,-13mm) -- (-123.5mm,-0mm) -- cycle;
                    \fill[color=corneryellow](0mm,-297mm) -- (0mm,-290.5mm) -- (-52.5mm,-290.5mm) -- (-61.7mm,-297mm) -- cycle;
                \end{tikzpicture}
            }
            \BgMaterial%
        }%
        \fi
}



%---------------
%   CHALLENGE PAGE BACKGROUND STYLE
%---------------

\newif\ifBgUseTwo % Lys's test for removing bg from pages

\backgroundsetup{scale=1,angle=0,opacity=1,contents={}}
\AddEverypageHook{%
    \ifBgUseTwo % Lys's test for removing bg from pages
    \backgroundsetup{contents={}
}%  Set an if/else conditional statement that checks whether the page is odd or not
    \ifthenelse{\isodd{\value{page}}}
        {% IF YES
            \SetBgPosition{current page.north west}
            \SetBgContents{
                \begin{tikzpicture}[remember picture,overlay,xshift=0mm,yshift=0mm]
                    \fill[color=chalgreen](0mm,-297mm) -- (0mm,-290.5mm) -- (95.8mm,-290.5mm) -- (105mm,-297mm) -- cycle;
                \end{tikzpicture}
            }
            \BgMaterial%
        }%
        {% IF NOT
            \SetBgPosition{current page.north east}
            \SetBgContents{
                \begin{tikzpicture}[remember picture,overlay,xshift=0mm,yshift=0mm]
                    \fill[color=chalgreen](0mm,0mm) -- (0mm,-13mm) -- (-105mm,-13mm) -- (-123.5mm,-0mm) -- cycle;
                    \fill[color=chalgreen](0mm,-297mm) -- (0mm,-290.5mm) -- (-52.5mm,-290.5mm) -- (-61.7mm,-297mm) -- cycle;
                \end{tikzpicture}
            }
            \BgMaterial%
        }%
        \fi
}


%------------------------------------------------------
%   HEADINGS STYLING
%------------------------------------------------------

\newcommand{\sectionnamesize}{\fontsize{40pt}{40pt}\selectfont}
\newcommand{\partsize}{\fontsize{32pt}{32pt}\selectfont} % Creates a command that sets font size to X
\newcommand{\chaptersize}{\fontsize{26pt}{26pt}\selectfont} % Creates a command that sets font size to Y
\newcommand{\headingonesize}{\fontsize{18pt}{18pt}\selectfont} % Creates a command that sets font size to 18
\newcommand{\subheadingsize}{\fontsize{15pt}{15pt}\selectfont} % Creates a command that sets font size to 15
\newcommand{\subsubheadingsize}{\fontsize{13pt}{13pt}\selectfont} % Creates a command that sets font size to 15

\newcommand{\segsize}{\fontsize{64pt}{64pt}\selectfont} % command for Segments title size
\newcommand{\issuesize}{\fontsize{20pt}{20pt}\selectfont}

%---------------
%   CHAPTER
%---------------

\makeatletter
  \renewcommand{\thechapter}{\two@digits{\value{chapter}}}
\makeatother

\makechapterstyle{box}{
  \renewcommand*{\printchaptername}{}
  \renewcommand*{\chapnumfont}{\headingfont\bfseries\fontsize{40pt}{40pt}\selectfont}
  \renewcommand*{\printchapternum}{
    \flushleft
    \begin{tikzpicture}
      \draw[fill,color=almostblack, xshift=-26mm] (0mm,0mm) rectangle (10.5mm, 26mm); % big box
      \draw[color=almostblack, xshift=-12mm] (10.5mm,13mm) node { \chapnumfont\thechapter }; % chapter number
      \draw[fill, color=almostblack, xshift=13mm] (0mm, 0mm) rectangle (0.8mm,26mm);   % separator
      %\draw[color=almostblack, xshift=26mm] (26mm,13mm) node {\printchaptertitle{test}};
    \end{tikzpicture}
  }
  \renewcommand*{\afterchapternum}{}
  \renewcommand\printchaptertitle[1]{\headingfont\chaptersize\bfseries%
    \hskip6.5mm \raisebox{10.5mm}{\parbox[s]{\textwidth-55mm}{##1}}
  }
}

\chapterstyle{box}
\newcommand{\byline}[1]{% command to add a by-line
\hskip15mm \raisebox{20mm}{\parbox{\textwidth-55mm}{\headingfont\subheadingsize by \textbf{#1}}}
\newline
}
\newcommand{\articlesub}[1]{% command to make article subtitles with the purple box (since subtitles don't seem to fit nicely with chapterstyles)
\begin{tikzpicture}
  \draw[fill,color=mainoffwhite] (0mm,0mm) rectangle (1mm,1mm);
  \draw[fill,color=purplevivid] (\textwidth+21mm,0mm) rectangle ++(-10.5mm, 13mm);
  \node[anchor=east] at (\textwidth+0mm, 5.5mm) {\headingfont\bfseries\headingonesize #1}; 
\end{tikzpicture}
}
\newcommand{\chalsub}[1]{% command for box color for challenges)
\begin{tikzpicture}
  \draw[fill,color=mainoffwhite] (0mm,0mm) rectangle (1mm,1mm);
  \draw[fill,color=watermelony] (\textwidth+21mm,0mm) rectangle ++(-10.5mm, 13mm);
  \node[anchor=east] at (\textwidth+0mm, 5.5mm) {\headingfont\bfseries\headingonesize #1}; 
\end{tikzpicture}
}



%---------------
%   SECTION
%---------------

\titleformat{\section} % we want to style the section headings only
  {\headingfont\bfseries\headingonesize}{\thesection}{1em}{}
  \titlespacing{\section}{0em}{0em}{0em}

%---------------
%   SUBSECTION
%---------------

\titleformat{\subsection} % we want to style the subsection headings only
  {\headingfont\bfseries\subheadingsize}{\thesection}{1em}{}
  \titlespacing{\subsection}{6.5mm}{0em}{0em}
 
\titleformat{\subsubsection} % we want to style the subsubsection headings only
  {\headingfont\bfseries\subsubheadingsize}{\thesection}{1em}{}
  \titlespacing{\subsubsection}{6.5mm}{0em}{0em}
  
%------------------------------------------------------
%   TEXT STYLING
%------------------------------------------------------

\setlength{\parindent}{3.2mm}   % sets the indentation of the first line of a paragraph
\setlength{\parskip}{12pt}      % sets the space between the paragraph and what's above it
\color{almostblack}

\newcommand{\highlight}[1]{{\color{purplevivid}{#1}}}

%------------------------------------------------------
%   ARTICLE-SPECIFIC STUFF
%------------------------------------------------------

% reduplication
\newbaarucmd[bits]{^}{\baarujuncture{\~{}}}

% infixes
\newbaarucmd[bits]{<}{\baaruopenbracket〈}
\newbaarucmd[bits]{>}{\baaruclosebracket〉}

% smoyding

\newcommand{\dex}[2]{\textbf{#1} \textit{`#2'}}
\def\baarusmoyd#1{(\textsc{5moyd} \##1)}
\newbaarulinetype{citation}{smoyd}{\baarusmoyd}

% ipa glossing
\def\baarubrack#1{[#1]}
\newbaarulinetype{unaligned}{ipa}{\baarubrack}

%------------------------------------------------------
%   LABELS & REF -- KUDOS TO AKAM
%------------------------------------------------------

% Wrap each article in \begin{article}{PREFIX}...\end{article}. PREFIX will
% be added to each cross-referencing label defined or used within the
% article. Two hooks, \PREFIXdefs and \endPREFIXdefs, are run at the start
% and end of the environment. The first can be used to load article-specific
% definitions or configuration, the second is probably useless.

% \makeatletter
% \AtBeginDocument{%
%   \let\seg@ref@saved\ref
%   \let\seg@label@saved\label
%   \let\seg@hyperref@saved\hyperref

%   \def\ref{\@ifstar\seg@ref@starred\seg@ref@}
%   \def\seg@ref@starred#1{\seg@ref@saved*{\segprefix#1}}
%   \def\seg@ref@#1{\seg@ref@saved{\segprefix#1}}

%   \def\label#1{\seg@label@saved{\segprefix#1}}

%   \def\hyperref{\@ifnextchar[\seg@hyperref@optarg\seg@hyperref@saved}%]
%   \def\seg@hyperref@optarg[#1]{\seg@hyperref@saved[\segprefix#1]}
% }

% \newenvironment{article}[1]{%
%   \ifcsdef{segprefix:#1}{%
%     \PackageError{segments}{Duplicate article prefix `#1'}\@eha
%   }{}%
%   \global\cslet{segprefix:#1}\relax
%   \def\segprefix{#1}%
%   \csuse{#1defs}%
% }{%
%   \csuse{end\segprefix defs}%
% }%
% \makeatother

\baabbrev[\FIRST]{1}{First person}
\baabbrev[\OBJ]{o}{Object}
\baabbrev[\SECOND]{2}{Second person}
\baabbrev[\SUBJ]{s}{Subject}
\baabbrev[\THIRD]{3}{Third person}
\baabbrev{1r}{Relating to first person}
\baabbrev{abl}{Ablative}
\baabbrev{abst}{Abstraction}
\baabbrev{absl}{Absolute}
\baabbrev{abs}{Absolutive}
\baabbrev{acc}{Accusative}
\baabbrev{act}{Active}
\baabbrev{adc}{Adjunct}
\baabbrev{adess}{Adessive}
\baabbrev{adjz}{Adjectivizer}
\baabbrev{adj}{Adjective}
\baabbrev{advz}{Adverbializer}
\baabbrev{adv}{Adverb}
\baabbrev{aff}{Affirmative}
\baabbrev{agn}{Agentive}
\baabbrev{all}{Allative}
\baabbrev{and}{Andative}
\baabbrev{an}{Animate}
\baabbrev{antess}{Antessive}
\baabbrev{aor}{Aorist}
\baabbrev{appl}{Applicative}
\baabbrev{apv}{Antipassive}
\baabbrev{atb}{Autobenefactive}
\baabbrev{atel}{Atelic}
\baabbrev{attr}{Attributive}
\baabbrev{aug}{Augmentative}
\baabbrev{aux}{Auxiliary}
\baabbrev{a}{Agent}
\baabbrev{ben}{Benefactive}
\baabbrev{caus}{Causative}
\baabbrev{cess}{Cessative}
\baabbrev{change}{Change of State}
\baabbrev{cl}{Classifier}
\baabbrev{cmpl}{Completive}
\baabbrev{cmpr}{Comparative}
\baabbrev{col}{Collective}
\baabbrev{com}{Comitative}
\baabbrev{cond}{Conditional}
\baabbrev{conj}{Conjunction}
\baabbrev{conn}{Connective particle}
\baabbrev{conneg}{Connegative}
\baabbrev{con}{Construct}
\baabbrev{cont}{Continous}
\baabbrev{coord}{Coordination}
\baabbrev{cop}{Copula}
\baabbrev{cp}{Complement phrase}
\baabbrev{cvb}{Converb}
\baabbrev{cyc}{Cyclical gender}
\baabbrev{c}{Complementizer}
\baabbrev{dat}{Dative}
\baabbrev{deadline}{Deadline}
\baabbrev{decl}{Declarative}
\baabbrev{def}{Definite}
\baabbrev{deic}{Deictic}
\baabbrev{del}{Delimitative}
\baabbrev{dem}{Demonstrative}
\baabbrev{dep}{Dependent}
\baabbrev{desid}{Desiderative}
\baabbrev{detr}{Detransitivizer}
\baabbrev{det}{Determiner}
\baabbrev{dim}{Diminutive}
\baabbrev{dir}{Direct}
\baabbrev{dis}{Distal/Distant}
\baabbrev{do}{Direct Object}
\baabbrev{dp}{Discourse particle}
\baabbrev{ds}{Different-subject}
\baabbrev{dum}{Dummy pronoun}
\baabbrev{dur}{Durative}
\baabbrev{du}{Dual}
\baabbrev{dvb}{Deverbal}
\baabbrev{ego}{Egophoric}
\baabbrev{emph}{Emphatic}
\baabbrev{erg}{Ergative}
\baabbrev{enc}{Enclitic}
\baabbrev{epis}{Epistemic}
\baabbrev{ess}{Essential} % if someone wants to make ESS essive than imo we should let them and get rid of this unless it comes up again
\baabbrev{ete}{Eternal gender}
\baabbrev{excl}{Exclusive}
\baabbrev{exist}{Existential}
\baabbrev{expl}{Expletive}
\baabbrev{exp}{Experiential; Direct Evidential}
\baabbrev{e}{Edible}
\baabbrev{fam}{Familiar}
\baabbrev{fin}{Finite Verb}
\baabbrev{fish}{Classifier for fish}
\baabbrev{foc}{Focus}
\baabbrev{fol}{Following stance}
\baabbrev{food}{Food}
\baabbrev{for}{Formal}
\baabbrev{frust}{Frustrative}
\baabbrev{fut}{Future}
\baabbrev{f}{Feminine}
\baabbrev{gen}{Genitive}
\baabbrev{ger}{Gerund}
\baabbrev{gno}{Gnomic}
\baabbrev{gnr}{General}
\baabbrev{hab}{Habitual}
\baabbrev{hon}{Honorific}
\baabbrev{hort}{Hortative}
\baabbrev{hyp}{Hypothetical Future}
\baabbrev{h}{Human}
\baabbrev{id}{Ideophone}
\baabbrev{imp}{Imperative}
\baabbrev{inal}{Inalienable Possession}
\baabbrev{inan}{Inanimate}
\baabbrev{inch}{Inchoative}
\baabbrev{incl}{Inclusive}
\baabbrev{inc}{Inceptive}
\baabbrev{indef}{Indefinite}
\baabbrev{ind}{Indicative}
\baabbrev{ine}{Inessive}
\baabbrev{infer}{Inferential}
\baabbrev{inf}{Infinitive}
\baabbrev{instr}{Instrumental}
\baabbrev{ins}{Instantive}
\baabbrev{int}{Intermediate tense}
\baabbrev{inv}{Inverse}
\baabbrev{ipfv}{Imperfective}
\baabbrev{irr}{Irrealis}
\baabbrev{it}{Itive}
\baabbrev{itr}{Iterative}
\baabbrev{itrt}{Intrative}
\baabbrev{juss}{Jussive}
\baabbrev{lat}{Lative}
\baabbrev{lcn}{Location Agreement}
\baabbrev{ld}{Leading Stance}
\baabbrev{lnk}{Linker}
\baabbrev{loc}{Locative}
\baabbrev{l}{Location Anaphor}
\baabbrev{mat}{Material}
\baabbrev{med}{Medial}
\baabbrev{mid}{Middle voice}
\baabbrev{min}{Minimal}
\baabbrev{m}{Masculine}
\baabbrev{name}{Personal Name}
\baabbrev{nec}{Necessitative}
\baabbrev{neg}{Negative}
\baabbrev{next}{Next}
\baabbrev{nfem}{Non-feminine}
\baabbrev{nfi}{Nonfuture imperfective}
\baabbrev{nfor}{Informal}
\baabbrev{nfut}{Non-Future}
\baabbrev{nha}{Non-human animate}
\baabbrev{nmlz}{Nominalizer}
\baabbrev{nml}{Animal Classifier}
\baabbrev{nom}{Nominative}
\baabbrev{npr}{Impersonal}
\baabbrev{npst}{Nonpast}
\baabbrev{nrz}{Non-realized}
\baabbrev{ntr}{Intransitive}
\baabbrev{nvl}{Nonvolitional}
\baabbrev{n}{Neuter}
\baabbrev{oblig}{Obligatory}
\baabbrev{obl}{Oblique}
\baabbrev{obv}{Obviative}
\baabbrev{opt}{Optative}
\baabbrev{org}{Origin}
\baabbrev{orn}{Ornative}
\baabbrev{pa}{Paucal}
\baabbrev{pcp}{Participle}
\baabbrev{perl}{Perlative}
\baabbrev{perm}{Permissive}
\baabbrev{per}{Personal}
\baabbrev{pfv}{Perfective}
\baabbrev{pf}{Patient Focus}
\baabbrev{place}{Place name}
\baabbrev{pl}{Plural}
\baabbrev{pn}{Pronoun, Proper noun}
\baabbrev{pol}{Polite}
\baabbrev{poss}{Possession}
\baabbrev{pos}{Possessive}
\baabbrev{pot}{Potential}
\baabbrev{pp}{Preposition Phrase}
\baabbrev{pret}{Preterite}
\baabbrev{prf}{Perfect}
\baabbrev{priv}{Privative}
\baabbrev{prl}{Prolative}
\baabbrev{prog}{Progressive}
\baabbrev{prop}{Proper article}
\baabbrev{prox}{Proximal}
\baabbrev{prsp}{Prospective}
\baabbrev{prs}{Present}
\baabbrev{pst}{Past}
\baabbrev{psv}{Passive}
\baabbrev{ptv}{Partitive}
\baabbrev{pt}{Plurative}
\baabbrev{punct}{Punctual}
\baabbrev{pv}{Preverb}
\baabbrev{p}{Patient}
\baabbrev{q}{Interrogative}
\baabbrev{quot}{Quotative}
\baabbrev{rdp}{Reduplication}
\baabbrev{rec}{Recent/Near tense}
\baabbrev{refl}{Reflexive}
\baabbrev{rel}{Relative}
\baabbrev{rem}{Remote tense}
\baabbrev{rep}{Reportative}
\baabbrev{res}{Resultative}
\baabbrev[\RR,sort=rr]{r/r}{Reflexive/reciprocal}
\baabbrev{rls}{Realis}
\baabbrev{rvs}{Reversative}
\baabbrev{rz}{Realized}
\baabbrev{sap}{Speech-Act-Participant}
\baabbrev{sbjv}{Subjunctive}
\baabbrev{sca}{Scalar/additive particle}
\baabbrev{sembl}{Semblative}
\baabbrev{seq}{Sequential}
\baabbrev{sg}{Singular}
\baabbrev{sim}{Simple aspect}
\baabbrev{ss}{Same-subject}
\baabbrev{stat}{Stative}
\baabbrev{st}{Singulative}
\baabbrev{sub}{Subordinator}
\baabbrev{supl}{Superlative}
\baabbrev{tel}{Telic}
\baabbrev{tem}{Temporary Gender} % gender is temporary; garlic bread is forever.
\baabbrev{top}{Topic}
\baabbrev{tri}{Trial}
\baabbrev{trn}{Transnumeral}
\baabbrev{trz}{Transitivizer}
\baabbrev{tr}{Transitive}
\baabbrev{trns}{Translative}
\baabbrev{vbz}{Verbializer}
\baabbrev{ven}{Venitive}
\baabbrev{via}{Vialis Argument}
\baabbrev{vid}{Verbal Identifier}
\baabbrev{vis}{Visual}
\baabbrev{voc}{Vocative}
\baabbrev{vol}{Volitional}
\baabbrev{vp}{Verb Phrase}
\baabbrev{v}{Verb}




\begin{document}

% 
\begin{center}
    \includegraphics[width=0.2\textwidth]{segmentslogo.png}
\end{center}

\pagestyle{empty}
\pagenumbering{gobble}

\section*{Instructions}

%Hi! This is the \LaTeX\hspace{.5mm} template for submitting articles for Segments, a Journal of Constructed Languages.

Please copy this project, and only edit the {\color{purplevivid} main.tex} document. Once you are done with your article, share the {\color{purplevivid} edit} link and send it to {\color{purplevivid} segments.journal@gmail.com}.

An article consists of the following, optional items in \highlight{purple}

\begin{enumerate}
    \item The title, subtitle, and author's name
    \item An introduction (250-800 characters incl. spaces)
    \item A phrase, or a handful of sentences, in the conlang, accompanied by its translation
    \item A body (two full pages would be a good minimum)
    \item {\color{purplevivid}A gloss of the phrase(s) from item 3}
    \item {\color{purplevivid}A few pictures, graphs, tables, or other useful graphics}
\end{enumerate}

Highlight important information using the command \textbackslash highlight: \highlight{highlighted text here}

\subsection*{Title, subtitle}

The title of your article will generally include the name of your conlang or project. The subtitle can be humorous, or descriptive, or somewhere in between!

Here are some examples of perfectly decent title -- subtitle combinations:

\begin{itemize}
    \item Mwaneḷe -- Ụṇḍẹṛḍọṭṣ
    \item ǂa ɳṵĩ -- A naturalistic click language
    \item Wistanian, a phonology of -- An odyssey into the sounds of Wistania
    \item An introduction to Valdean -- A written language
\end{itemize}

\subsection*{Graphics, tables}

The inclusion of graphics is there so you can indicate where the images you want to include in your article are best placed, according to you. We will try our best to accomodate your choices, however please be aware that differences may arise in the final product to facilitate laying it out. We will keep these differences as small as possible, and will send you a copy of your article for review before publication.

You can include them with the command \verb|\includegraphics[]{name_of_image}| (see example above in this document's code with the logo).

\subsection*{Notes, comments, and asides}

You can add notes and asides as footnotes to this document. We will integrate them into the article.\footnote{Use \textbackslash footnote\{\} to include a footnote!}




\newpage




\section*{Styling guide: Format \& examples}

You can copy code from this section and alter it to fit your needs.

\subsection*{Paragraph styles}

\subsubsection*{Conlang, translations}

Words, or morphemes, of the conlang should be \textbf{bolded} if they are inside of an English sentence. 

English translations should be in \textit{italics}, and

\begin{itemize}
    \item surrounded by \highlight{`single quotemarks'} if it is a single word, or a definition;
    \item surrounded by \highlight{``double quotemarks"} inside examples and glosses.
\end{itemize}

IPA transcriptions should be left in plain, normal text, unless specific phonemes are talked about within an English sentence, in which case they must be in \textit{italics}, and can optionally be \highlight{highlighted in purple}.

Here is an example: \\
The verb \textbf{puqwe} /pɯqʷɛ/ \textit{`to clean up, to tidy up, to arrange, to put in order'} is used with bananas. here, the phoneme \highlight{\textit{qʷ}} becomes [t͡l] through assimilation of the PoA of /j/ which is not present in this word at all, and this sentence makes no sense.

\subsubsection*{Conscripts}

The best way to include a conscript is to send, along with your article, a high resolution .png, or a .svg of it, along with all relevant instructions for its placement.

\subsubsection{Custom fonts}

If your conlang makes use of a writing system for which several fonts are available, please provide us with a .ttf or .otf file of the one you would like us to use, \textbf{after making sure that font is under a license that would allow for a third party publication to make use of it.}

Your TeX file should come with the necessary packages and commands already implemented at the top of \textbf{main.tex}, right after line 3, which inputs our \textbf{specifications.tex} file.

If no such font is available, a high resolution .png, or a .svg of it would be ideal.

\subsubsection{Additional specifications}

Latin locutions, such as \textit{etc}, \textit{e.g.}, \textit{i.e.} and others should be in \textit{italics}.

\subsection*{Tables}

Headers for both columns and rows are in bold. Row headers are aligned to the left, and everything else to the center.

The only two lines drawn are right next to the headers, and run across the entirety of the table.

\begin{table}[H]
	\centering
	\begin{tabu}{$>{\bfseries}l|^c^c^c}     % ^ before a column whose header should be bold
	\rowstyle{\bfseries}                    % The next row will be bold
		& Labial & Alveolar & Dorsal \\
		\hline
		Nasal       & m     		   & n        		   & q      \\
		Stop        & b     		   & t d      		   & k g    \\
		Fricative   & f     		   & s x     		   & c h    \\
		Trill       &        		   & r       		   &        \\
		Approximant & w    		   & l      		   & j      
	\end{tabu}
	\caption{A single table}
	\label{cons-inv}
\end{table}

\begin{figure}[H]

\begin{multicols}{2}
\centering

    \begin{tabu}{$>{\bfseries}l|^c^c}
	\rowstyle{\bfseries}
    & Front & Back \\
    \hline
    High & i & u \\
    Mid & ɛ & ɔ\\
    Low & a & ɒ
    \end{tabu}

    \begin{tabu}{$>{\bfseries}l|^c^c}
	\rowstyle{\bfseries}
    & Unrounded & Rounded \\
    \hline
    High & i & u \\
    Mid & ɛ & ɔ\\
    Low & a & ɒ
    \end{tabu}

\end{multicols}

\caption{Here are two tables, on the same level}
\end{figure}
\clearpage 
\subsection*{Vowel charts}

Vowel charts can be constructed using the \highlight{\texttt{tikzvowel}} package. To create a vowel chart, begin a \texttt{vowel} environment. Here is an example vowel chart. You can modify the size of the trapezoid by adjusting the \textbf{scale} parameter.

\begin{figure}[H]
  \centering
  \begin{vowel}[scale=0.75]
    \vpoint{0}{0}{i}
    \vpoint{0}{2}{u}
    \vpoint{1.5}{1}{ə}
    \vpoint{3}{1}{a}
    \vblob{-0.1, 2.1}{1.2, 2.1}{1.2, 1.7}{-0.1, 1.7}
    \varrow{a}{i}
  \end{vowel}
  \caption{Phonemic Vowel Inventory}
  \label{table:vowel_phonemes}
\end{figure}
\begin{verbbox}
  \begin{vowel}[scale=0.75]
    \vpoint{0}{0}{i}
    \vpoint{0}{2}{u}
    \vpoint{1.5}{1}{ə}
    \vpoint{3}{1}{a}
    \vblob{-0.1, 2.1}{1.2, 2.1}{1.2, 1.7}{-0.1, 1.7}
    \varrow{a}{i}
  \end{vowel}
\end{verbbox}
\fbox{\theverbbox} 

Add vowels using the command \texttt{\textbackslash vpoint\{y\}\{x\}\{V\}}, where \texttt{x} is the horizontal position of the vowel, \texttt{y} is the vertical position of the vowel, and \texttt{V} is the vowel letter you want to add. The origin is set at the top left corner. If \textbf{scale} is set to 1, then each horizontal or vertical line is one unit. For example in figure \ref{table:vowel_phonemes}, \textbf{a} is inserted three columns down and one row to the right, using the command \texttt{\textbackslash vpoint\{3\}\{1\}\{a\}}. The x and y positions don't have to be integers: to place the schwa between the mid-closed line and the mid-open line, we used y=1.5.

In addition to adding vowel letters, the \texttt{tikzvowel} package can add arrows to show diphthongs or vowel shifts, as well as outlines to show ranges within the vowel chart. Arrows can be added using the command \texttt{\textbackslash varrow\{a\}\{b\}}, which draws an arrow from vowel a to vowel b. Here, the arrow from \textbf{a} to \textbf{i} is drawn with \texttt{\textbackslash varrow\{a\}\{i\}}. 

Dotted outlines can be drawn with the \texttt{\textbackslash vblob} command, which takes a series of coordinates and draws a curved shape with those coordinates as rough corners. The dotted blob showing the vowel space occupied by the phoneme /u/ in figure \ref{table:vowel_phonemes} is drawn with the command \texttt{\textbackslash vblob\{-0.1, 2.1\}\{1.2, 2.1\}\{1.2, 1.7\}\{-0.1, 1.7\}}. 

Beware that this package gives coordinates in (y, x) and uses the top left as the origin point (unlike what they tend to teach in school!). Other than that point, this is an easy and versatile tool to draw vowel diagrams!

\newpage
\section*{Using Baarux}
\counterwithout{exesi}{section} % needed to get baarux to number the way I want outside of chapters

In Segments, we use the package \highlight{\texttt{baarux}} by community member \highlight{Akam Chinjir} to typeset examples. The current version as of writing is \highlight{0.9.9}, which we've included in this template. As Akam pushes new versions, we'll update our documents including this guide. \texttt{baarux} is a very powerful package, so we won't go into everything it can do, but here's an introduction to its functionalities. % link to akam's documentation as extra reference when it's done

The core of \texttt{baarux} is the \highlight{examples} environment, which creates numbered examples. Declare a new example within the examples environment using \texttt{\textbackslash ex}. You can nest examples environments to make subexamples (up to four levels deep, but if you need more than that, then stop and think about your life choices). You can use the \texttt{\textbackslash label} and \texttt{\textbackslash ref} commands to label examples and refer back to them by number, like example \ref{ex:protest}.

I'll follow each example with the code used to make it, so that you can get a sense of how \texttt{baarux} is used.

\begin{examples}
    \ex This is a test example
    \ex And these are some test subexamples:
    \begin{examples}
        \ex Test
        \ex Test
        \ex Protest \label{ex:protest}
    \end{examples}
\end{examples}
\begin{verbbox}
\begin{examples}
    \ex This is a test example
    \ex And these are some test subexamples:
    \begin{examples}
        \ex Test
        \ex Test
        \ex Protest \label{ex:protest}
    \end{examples}
\end{examples}
\end{verbbox}
\fbox{\theverbbox} 

Within the examples environment, \texttt{baarux} defines several line types. Here are the predefined plain line types:
\begin{itemize}
    \item \highlight{\texttt{\textbackslash preamble}} gives an unformatted line
    \item \highlight{\texttt{\textbackslash script}} is bolded and is meant for the in-language text of the example
    \item \highlight{\texttt{\textbackslash tr}} is automatically surrounded by quotation marks and is meant for the translation
    \item \highlight{\texttt{\textbackslash context}} prints the line preceded by `Context: '
    \item \highlight{\texttt{\textbackslash alt}} prints the line inside quotation marks and preceded by \textit{`Or: '} and is used to introduce alternate translations
    \item \highlight{\texttt{\textbackslash intended}} prints the line inside quotation marks and preceded by \textit{`Intended: '} and can be used to give the intended meaning of an ungrammatical or questionable sentence 
    \item \highlight{\texttt{\textbackslash not}} prints the line inside quotation marks and preceded by \textit{`Not: '} and can be used to give examples of what a sentence \emph{doesn't} mean to contrast with what it does mean
\end{itemize} 
There are also two predefined line types that give right-aligned comments: \highlight{\texttt{\textbackslash lect}} prints a plain right-aligned comment and is meant to give the name of the language or variety used in the example, while \highlight{\texttt{\textbackslash source}} prints a right-aligned comment in parentheses and is meant to show where the text in the example is coming from. (These are the intended uses, but of course you can be creative! Reach out if you're interested in learning how to define your own line types.)

Example \ref{ex:seoina} shows a sentence written in the conlang Seoina taken from the \href{https://www.reddit.com/r/conlangs/comments/haaf0c/1278th_just_used_5_minutes_of_your_day/} {1278\textsuperscript{th} Just Used 5 Minutes of your Day Challenge} from r/conlangs. 

\begin{examples}
    \ex \label{ex:seoina}
    \lect Seoina
    \script Sa kel si la deol, aloi la kipia nolra peu deoi.
    \tr How you treat yourself, that's how people will think to treat you.
    \alt People will think to treat you however you treat yourself.
    \source \textsc{5moyd} \#1278
\end{examples}

\begin{verbbox}
\begin{examples}
    \ex \label{ex:seoina}
    \lect Seoina
    \script Sa kel si la deol, aloi la kipia nolra peu deoi.
    \tr How you treat yourself, that's how people will think to treat you.
    \alt People will think to treat you however you treat yourself.
    \source \textsc{5moyd} \#1278
\end{examples}
\end{verbbox}
\fbox{\theverbbox}

You can also easily do glossed example sentences with \texttt{baarux}. In addition to plain lines and right-aligned comments, the package defines three types of lines that automatically align for glossing:

\begin{itemize}
    \item \highlight{\texttt{\textbackslash words}} is automatically bolded and meant for words in the language
    \item \highlight{\texttt{\textbackslash bits}} is plain and meant for words, morphemes, affixes, clitics...you know, `bits'
    \item \highlight{\texttt{\textbackslash gloss}} is italicized except for glossing abbreviations, which are automatically detected and rendered in upright smallcaps--more on that later
\end{itemize}

The \texttt{baarux} package comes with a companion \texttt{baabbrevs} which is used to define glossing abbreviations. Abbreviations are defined using the command \texttt{\textbackslash baabbrev}, which takes two arguments: the abbreviation in all lowercase followed optionally by the full term that the abbreviation is short for. You can see examples of this in the \highlight{abbreviations.tex} file under the Required folder here. Please use glossing terms that are already present in that list whenever possible; if you have additional glossing abbreviations needed that are \textbf{not} in that list, please add them to the top of your \highlight{main.tex} file so they can be properly added when Segments is compiled completely.

To make a glossed example, add your words or bits, separated by hyphens or equals signs wherever you want there to be junctures. Then add a gloss row with \textbf{only spaces} between the bits of gloss, and with defined abbreviations in all caps. \texttt{baarux} will automatically break them up and align them, add hyphens and equals signs to the gloss row corresponding to whatever's in the bits row, and turn abbreviations into \textsc{small caps}. You can use parentheses, colons and periods in the gloss line without messing things up. You can change the alignment of individual bits by using spaces to separate juncture markers in the bits line.

\begin{examples}
    \ex \label{ex:mwanele}
    \lect Mwaneḷe
    \script Mwana xalolo xo tetesi exeŋi ekwulife.
    \bits Mwana xalolo =xo ta- e- tesi e- xeŋi e- kwu- life
    \gloss NAME fear:NFI DP C APV exceed APV be.below APV VEN arrive
    \tr Mwana is afraid she arrived too late.
    \source \textsc{5moyd} \#1270
\end{examples}

\begin{verbbox}
\begin{examples}
    \ex \label{ex:mwanele}
    \lect Mwaneḷe
    \script Mwana xalolo xo tetesi exeŋi ekwulife.
    \bits Mwana xalolo =xo ta- e- tesi e- xeŋi e- kwu- life
    \gloss NAME fear:NFI DP CMP APV exceed APV be.below APV VEN arrive
    \tr Mwana is afraid she arrived too late.
    \source \textsc{5moyd} \#1270
\end{examples}
\end{verbbox}
\fbox{\theverbbox}

Sometimes it can be useful to have items in glossed lines that extend across multiple columns. You can do this easily with the \highlight{\texttt{\textbackslash MC}} command. If you put \texttt{\textbackslash MC} followed by a number \textit{n}, that tells \texttt{baarux} that the next word should occupy \textit{n} cells. 

On the other hand, sometimes you want multiple words in a single cell or column. If you want to include a string with a space in it, then you can surround it by brackets \{like this\} to tell \texttt{baarux} to treat it as one unit, instead of treating the space as a word/morpheme break.

\begin{examples}
    \ex
    \lect 3eyri
    \words Čelt \MC4 tanolhka \MC2 piko mi \MC3 tagökalh?
    \bits čelt t -ano  -lh =ga pik =o mi t -gök -lh
    \gloss man 3ACC sleep 3NOM in boy NOM what 3ACC eat 3NOM
    \gloss man \MC4 {while he slept} \MC2 {the boy} what \MC3 {he ate it}
    \tr What did the boy eat while the man was sleeping?
    \source \textsc{5moyd} \#1179
\end{examples}

\begin{verbbox}
\begin{examples}
    \ex
    \lect 3eyri
    \words Čelt \MC4 tanolhka \MC2 piko mi \MC3 tagökalh?
    \bits čelt t -ano  -lh =ga pik =o mi t -gök -lh
    \gloss man 3ACC sleep 3NOM in boy NOM what 3ACC eat 3NOM
    \gloss man \MC4 {while he slept} \MC2 {the boy} what \MC3 {he ate it}
    \tr What did the boy eat while the man was sleeping?
    \source \textsc{5moyd} \#1179
\end{examples}
\end{verbbox}
\fbox{\theverbbox}

You can already see that \texttt{baarux} is a very powerful package, and this is just the tip of the iceberg! If there's something you want to do with your examples, take a look at Akam's documentation or ask around. Chances are you can find a way.


\newpage
\mbox{}
\newpage

\pagenumbering{arabic}
\pagestyle{fancy} % Comment this to remove the help pages

%------------------------------------------------------
%   ↓   INTRODUCTION    ↓
%------------------------------------------------------

\chapter[Randomizing Language Features][Tamgwa]{Randomizing Language Features in Tamgwa}

\byline{armytag} % add your name as you wish to be credited
\articlesub{When and how to roll the dice} % use this command to insert a subtitle

\thispagestyle{fancy}
\BgUsetrue

%------------------------------------------------------
%   ↓   ARTICLE BODY    ↓
%------------------------------------------------------

    This article outlines how randomization was used to generate the initial features of \textbf{Tamgwa}, and the process of making those random features work coherently together.

\section*{What is randomization?} % * for unnumbered section

    Randomization, as used in this article, refers to a process in which a set of possibilities are narrowed down in a random manner.  Examples of this might be creating a list of choices and rolling some dice to select from it, or writing options on a deck of flash cards to draw from.  Digital tools can often help with randomization, especially for conditional or rule-based features like syllable and word generation which depend on a language's phonotactics.  Artificial Intelligence tools like ChatGPT could also be considered randomization tools under this definition, but they were not used in the randomization process of \textbf{Tamgwa} outlined in this article so they have been ignored here; AI tools are certainly not needed for randomization, as there are a variety of simpler techniques available.

    Randomization offers several benefits when it comes to language creation.  It works best as a tool in a language creator's toolbelt, augmenting their creation process without dominating it.  This can happen in several ways.  First, randomization can help a language creator generate content like phonemes, words, grammar structures, \textit{etc}.  This benefit should be obvious to anyone who winces remembering the half-finished dictionary they haven't touched in months. Second, randomization can counter a conlanger's linguistic biases.  Few of us are familiar with more than a handful of languages, and that limited exposure can likewise limit the form our conlangs take.  If you can't help but make Englishy conlangs, chance can steer you elsewhere.  Lastly, randomization can help overcome the indecision we sometimes feel in the face of endless linguistic possibility.  There are so many ways that human language conveys ideas, and sometimes it can be almost impossible to wade through the myriad options to settle on the \textbf{one} this language will use.

\section*{Laying the foundation} % * for unnumbered section

    To exemplify the randomization process, I will walk through the method I used for starting \textbf{Tamgwa} with almost entirely randomized features.  To reiterate the definition above, randomization here will always involve first identifying a set of possible options, then randomly selecting from that set.  For creating a list of possibilities I relied heavily on the invaluable World Atlas of Langauge Structures (WALS).\footnote{https://wals.info}  WALS has a collection of chapters and associated maps showing the distribution of a variety of language features, from phonology to morphology to syntax and more.  Each chapter lists the various realizations of that language feature along with the number of languages using each variation.  These numbers give an idea of how widespread each variation is and allow us to weight the randomization in a more naturalistic way.  In other words, we can use the chapters in WALS to randomize a set of language features that mimics the tendencies of natural languages.

     A typical randomization procedure can been seen in how noun plurality in \textbf{Tamgwa} was randomized using WALS chapters 33 \& 34.  Chapter 33 covers how plurality is encoded, such as using affixes like English's \textit{-s}.  The WALS authors list the following possibilities along with the number of languages which employ each option.

\begin{table}[H]
	\centering
	% \begin{tabu}{$>{\bfseries}l|^c^c^c^c^c}
	\begin{tabu}{>{\bfseries}l|^c^c}
	\rowstyle{\bfseries}
        Language Feature & Language Count & Randomizer Range \\
		\hline
        Plural prefix & 126 & 1 - 126 \\
        Plural suffix & 513 & 127 - 639 \\
        Plural stem change & 6 & 640 - 645 \\
        Plural tone & 4 & 646 - 649 \\
        Plural by complete stem reduplication & 8 & 650 - 657 \\
        Morphological, but no method primary & 60 & 658 - 717 \\
        Plural word & 170 & 718 - 887 \\
        Plural clitic & 81 & 888 - 968 \\
        No plural & 98 & 969 - 1066\\
	\end{tabu}
	\caption{WALS Chapter 33: Coding of Nominal Plurality}
	\label{wals-33}
\end{table}

    I have added a Randomizer Range column, which simply shows how a random number should map back to the feature options.\footnote{To create the Randomizer Ranges yourself, simply take the high end of the previous range, add 1 to it to get the low end of the new range, and add the number of languages to it to get the high end of the new range}  In this case a number can be randomly generated between 1 and 1066 (inclusive) to see which range it falls within.  For \textbf{Tamgwa} the random number was 814, so \textbf{Tamgwa} will encode plurality with an independant word rather than some kind of morphological inflection.  The same process can be used with Chapter 34, \textit{Occurrence of Nominal Plurality} to randomize when the plural word should be used.  The random number for this feature was ultimately 118 which corresponds to the option "Plural in all nouns, always optional", thus all nouns in \textbf{Tamgwa} can be optionally marked for plurality by adding the plural word somewhere in the sentence.  It is worth noting that sometimes there can be conflict between randomized features, such as if chapter 34 had been randomized to "No nominal plural" despite the fact that \textbf{Tamgwa} already has a plural word based on the feature in chapter 33.  In that case one could either exclude the conflicting option beforehand, or re-randomize if it is selected.

    When first creating \textbf{Tamgwa} I used the process above to randomize features for most of the WALS chapters.  This provided a wide grammatical basis upon which to build the rest of the language.  Some of the features naturally contradicted one another, as one would expect from a completely random result, but not all of those contradictions proved to be intractable.  The next section describes an interesting feature that developed in order to resolve one such conflict.

\section*{Personal disagreement} % * for unnumbered section

    The randomized features of \textbf{Tamgwa} are (surprisingly) mostly consistent, giving shape to an analytic language with very little inflectional morphology.  There are still some inconsistencies however, the most glaring of which was between the features for chapters 26 and 102.  Chapter 26, \textit{Prefixing vs. Suffixing in Inflectional Morphology}, uses a so-called \textit{affix index} to measure how much a language uses affixes.  The authors list all the features which contribute to calculating the \textit{affix index} in the chapter's description, but what is relevant here is that affixes marking a verb's agent (A) add 2 to the \textit{affix index} and affixes marking its patient (P) add 1.  If the total \textit{affix index} is 2 or less then the language is categorized under "little to no inflectional morphology", which is the case for \textbf{Tamgwa}.  However, the randomized result for chapter 102, \textit{Verbal Person Marking}, specifies that verbs in \textbf{Tamgwa} must have "person marking of both A and P arguments" which would seem to necessitate an affixing index of at least 3 and thus contradict the value generated for chapter 26.

    Although this contradiction seems to require one or both of the features to be re-randomized, there are some details in the WALS definitions that can actually help resolve the contradiction.  In chapter 26, the authors clarify that person-marking clitics that can also attach to non-verbs are not included as part of the \textit{affix index} calculation.  On the other hand, chapter 102 does count clitics in its analysis as long as the clitics can sometimes attach to the verb.  Therefore, if at least one of the person markers is a clitic that sometimes attaches to verbs and sometimes to non-verbs, then the \textit{affix index} can be kept to 2 or less while still marking the verb for both A and P.

    An real-world example of a language with verbal person marking that uses affixes and clitics together is the Sorani dialect of Kurdish.  According to the analysis by Gharib and Pye\footnote{Gharib, H., \& Pye, C. (2023) The clitic status of person markers in Sorani Kurdish. Kansas Working Papers in Linguistics, 39, 57-65. https://doi.org/10.17161/1808.27692}, one of the verb arguments is marked by an affix on the verb and the other argument is marked by a clitic on either the verb or a preceding adverb.  Whether the clitic is marking the agent (A) or the patient (P) of the action depends on whether it is using the Perfective or Imperfective aspect, as seen in the example below.

\begin{examples}
    \ex
    \lect Sorani Kurdish
    \words \MC2 dat \MC2 benim
    \bits da =t beni -m
    \gloss IPFV OBL.2SG see.IPFV 1SG
    \tr I see you.
    \ex
    \words \MC3 benemeet
    \bits beni =m -eet
    \gloss see.PFV OBL.1SG 2SG
    \tr I saw you.
    \ex
    \words \MC3 benetm
    \bits bene =t -m
    \gloss see.PFV OBL.2SG 1SG
    \tr You saw me.
    \source Gharib \& Pye 2023
\end{examples}

    In Sorani Kurdish the oblique clitic is used to mark the patient of imperfective verbs and the agent of perfective verbs.  \textbf{Tamgwa} can use a simpler system, where the clitic marks the same verbal argument and attaches to the first verb/adverb in the verb phrase regardless of aspect.  Thus for our purposes, it is sufficient to see that examples (1) and (3) show the OBL.2SG clitic attaching to the aspect marker if it is present, and the verb if it is not.

\section*{Phonology}

    Based on the person-marking system outlined above it is possible to start generating a full set of pronouns, affixes, and clitics, but first the phonology of \textbf{Tamgwa} must be established.  WALS has nineteen chapters dedicated to phonology, all of which have been randomized in \textbf{Tamgwa}.  Here are some of the more significant results:

\begin{itemize}
    \item Moderately large consonant inventory (26-33 consonants)
    \item Small vowel inventory (2-4 vowels)
    \item Voicing distinction in plosives but not fricatives
    \item Missing some of /p t k b d g/ but not exactly /p/, /g/, or both
    \item Has all common consonants (bilabials, fricatives, and nasals)
    \item Has no uncommon consonants (clicks, labial-velars, pharyngeals, or 'th' sounds)
    \item Has no uvulars
    \item Has no glottals
    \item Has lateral /l/ as well as obstruent laterals
    \item Has velar nasal, including word-initially
    \item No contrastive vowel nasalization
    \item No front rounded vowels
    \item Moderately complex syllable structure (up to CCVC)
\end{itemize}

    Although it would be feasible to further randomize the phonology by generating random phonemes until all the features are met, I preferred to choose the phonemes manually.  Given the large consonant inventory and small vowel inventory, I decided to take inspiration from the 2-vowel analysis of Mandarin phonology.  This means \textbf{Tamgwa} will have two underlying vowels /a/ and /ə/, but a set of labialized and palatalized consonants will trigger /ə/ to be realized as /u/ or /i/ respectively while losing their secondary articulation.  For examples the syllables /tʷə/ and /tʲə/ would become /tu/ and /ti/, but /ta/, /tʷa/, /tʲa/, and /tə/ would remain unchanged.  From here more consonants were added to fill out the rest of the consonant requirements, ultimately giving the table below.

\begin{table}[H]
	\centering
	% \begin{tabu}{$>{\bfseries}l|^c^c^c}
	\begin{tabu}{>{\bfseries}l|^c^c^c}
	\rowstyle{\bfseries}
        & Labial & Alveolar & Velar \\
		\hline
        Nassals              & m̥ m   & n̥ n & ŋ̊ ŋ   \\
        Plosives             & p b   & t   & k g   \\
        Labialized Plosives  & pʷ bʷ & tʷ  & kʷ gʷ \\
        Palatalized Plosives & pʲ bʲ & tʲ  & kʲ gʲ \\
        Fricatives           & f     & s   & x     \\
        Trill                &       & r   &       \\
        Laterals             &       & ɬ l &       \\
	\end{tabu}
	\caption{Consonants in Tamgwa}
	\label{consonants}
\end{table}

    As for vowels, \textbf{Tamgwa} can be analyzed as having two underlying vowels or four surface vowels, but either analysis falls within the feature randomized before.  Here is a table with the four-vowel analysis.

\begin{table}[H]
	\centering
	% \begin{tabu}{$>{\bfseries}l|^c^c^c}
	\begin{tabu}{>{\bfseries}l|^c^c^c}
	\rowstyle{\bfseries}
        & Front & Central & Back \\
		\hline
        High & i &   & u \\
        Mid  &   & ə &   \\
        Low  &   & a &   \\
	\end{tabu}
	\caption{Vowels in Tamgwa}
	\label{vowels}
\end{table}

\section*{Producing Pronouns}

    The features for \textbf{Tamgwa} were randomized such that pronouns will use words with distinct person-number stems and a clusivity distinction in 1PL pronouns.  The pronouns will have nominative and accusative forms in line with the separate A and P person-markers, as well as a form to indicate possession; the forms will be called nominative, accusative, and genetive forms respectively (although there is no case system for general nouns).  Altogether this entails 7 person-number combinations and 5 forms (3 cases, 1 affix, and 1 clitic form) for a total of 35 distinct "words".

    There are plenty of ways to generate a list of words and syllables, including digital tools such as Vulgarlang\footnote{https://www.vulgarlang.com}, but for \textbf{Tamgwa} I have created a Python script that is tailored to its phonology and phonotactics.  This script was used to generate 35 words that were randomly assigned to the different pronoun forms needed, as shown in the table below.

\begin{table}[H]
	\centering
	% \begin{tabu}{$>{\bfseries}l|^c^c^c^c^c}
	\begin{tabu}{>{\bfseries}l|^c^c^c^c^c}
	\rowstyle{\bfseries}
        & NOM & Affix & ACC & Clitic & GEN \\
		\hline
        1SG     & hnja      & ki    & gya       & hnja  & hma    \\
        2SG     & lam       & man   & ba        & be    & pu     \\
        3SG     & le        & na    & bwa       & nji   & pe     \\
        1PL.INC & ga        & fe    & hmuhl     & pix   & re     \\
        1PL.EXC & pwa       & tenj  & hla       & le    & pu     \\
        2PL     & bihl      & tu    & bya       & bes   & hnjanj \\
        3PL     & xa        & ti    & pwax      & njar  & bwa    \\
	\end{tabu}
	\caption{Random pronoun table}
	\label{rand-pronoun}
\end{table}

    Although there are some similar (even identical) words in the table, there is no immediate pattern.  These words were manually rearranged after the randomization to add more consistency. For example, all of the words beginning with \textbf{hnj} or \textbf{nj} were moved into the same row, or words with closed syllables were moved to the GEN column.  The final rearranged table for \textbf{Tamgwa} is printed below, with obligatory plural marker \textbf{tyan}\footnote{\textit{N.B.}\textbf{tyan} is an optional plural marker for all nouns, but is obligatory for plural pronouns.  The feature was randomly generated for WALS chapter 35.} added to NOM and ACC forms of plural pronouns.

\begin{table}[H]
	\centering
	% \begin{tabu}{$>{\bfseries}l|^c^c^c^c^c}
	\begin{tabu}{>{\bfseries}l|^c^c^c^c^c}
	\rowstyle{\bfseries}
        & NOM & Affix & ACC & Clitic & GEN \\
		\hline
        1SG     & hnja      & -nji  & hnja      & =njar & hnjanj \\
        2SG     & fe        & -pe   & pwax      & =pwa  & pix    \\
        3SG     & bes       & -hla  & hmuhl     & =hma  & man    \\
        1PL.INC & tu tyan   & -ti   & re tyan   & =na   & tenj   \\
        1PL.EXC & xa tyan   & -ki   & gya tyan  & =ga   & bya    \\
        2PL     & bwa tyan  & -be   & bwa tyan  & =ba   & bihl   \\
        3PL     & pu tyan   & -le   & pu tyan   & =le   & lam    \\
	\end{tabu}
	\caption{Rearranged pronoun table}
	\label{arr-pronoun}
\end{table}

\section*{Revealing Results}

    In addition to these pronominal forms, a set of preverbal tense-aspects words are required for the clitics to attach to in order to fully illustrate \textbf{Tamgwa}'s new person-marking syntax.  Once again pulling from real-world examples, \textbf{Tamgwa} will mirror the set of aspect distinctions described in Timothy Ajani's 2001 paper on Yoruba.\footnote{Ajani, T. 2001. \textit{Aspect in Yoruba and Nigerian English}}  According to this analysis, we will need 4 aspect words marking Realis Imperfective, Irrealis Perfective, Irrealis Imperfective, and Relational aspects.  The Realis Perfective is unmarked and thus assumed when none of these separate aspect words occur.  The aspect words can also occur together to create more nuanced aspectual distinctions, such as Realis Imperfective and Irrealis Imperfective denoting a Habitual aspect when used together.  For more details on this system, please refer to Ajani's paper.

    The same Python script used to generate the pronouns can be used to generate these 4 aspect words.  However, because the number of words is relatively small, the script can be set to generate 2-3 times as many words and the final choices can be manually selected.  The results are listed in the table below.  

\begin{table}[H]
	\centering
	% \begin{tabu}{$>{\bfseries}l|^c^c^c^c^c}
	\begin{tabu}{>{\bfseries}l|^c}
	\rowstyle{\bfseries}
        Aspect & Word \\
		\hline
        RLS.PFV  & <none> \\
        RLS.IPFV & kwa    \\
        IRR.PFV  & hna    \\
        IRR.IPFV & pi     \\
        REL      & gya    \\
	\end{tabu}
	\caption{Aspect words in Tamgwa}
	\label{aspect-words}
\end{table}

    These aspect words can now be used to test the person-marking clitic behaviour in some example sentences below.  Remember that the clitics will attach to the end of the first word in the aspect-verb structure.

\begin{examples}
    \ex
    \lect Tamgwa
    \words \MC3 kyamunjipwa
    \bits kyamu -nji =pwa
    \gloss see 1SG.NOM 2SG.ACC
    \tr I saw you.
    \ex
    \words \MC2 kwapwa \MC2 kyamunji
    \bits kwa =pwa kyamu -nji
    \gloss RLS.IPFV 2SG.ACC see 1SG.NOM
    \tr I see you.
    \ex
    \words \MC2 pipwa \MC2 kyamunji
    \bits pi =pwa kyamu -nji
    \gloss IRR.IPFV 2SG.ACC see 1SG.NOM
    \tr I will see you.
    \ex
    \words \MC2 pipwa kwa \MC2 kyamunji
    \bits  pi =pwa kwa kyamu -nji
    \gloss IRR.IPFV 2SG.ACC RLS.IPFV see 1SG.NOM
    \tr I generally (habitually) see you.
\end{examples}

\section*{Conclusion} % * for unnumbered section

    Hopefully this article has given you some ideas about how to use randomization to augment your conlanging process.  You can randomize everything from syntax to morphology to phonology to lexicon and more.  This is particularly useful if you're struggling with choice paralysis or not sure where to start with certain part of your language, and you can even go as far as bootstrapping most of a language's features through random generation.  You can also go about doing your randomization in whatever way you'd like.  You write down some options on flash cards to randomly draw from.  However you choose to use randomization, if at all, I wish you luck and can't wait to see what you come up with.

\end{document}
